% 本文件由账本已入列事件生成；锚点首次呈现仅连接可核的年序与材料限度。
\section{元朝区域治理与终局的加密年序}
以下补入按同年并列的政权与地区组织，而不是以一个朝廷的年号吞没其对手。每个可点击事件名保留账本 ID；叙述只将材料明确记录的动作置入年序，并把实施、规模、损失与动机留在可检验的限度内。

\subsection{1272年：元、元与南宋}
《元史》为这一年留下了可回查的年序入口。

\noindent\textbf{元。}\eventanchor{SYD-1272-YUAN-001}{:今遣汝等，水陸並進，布告遐邇，使咸知之}；\eventanchor{SYD-1272-YUAN-145}{敕杖合丹，斥無入宿衞，謫往西川效死軍中，餘定罪有差}；\eventanchor{SYD-1272-YUAN-217}{從其請，仍降璽書，遣使諭江陵府制置司及高達已下官吏軍民}；\eventanchor{SYD-1272-YUAN-289}{帝是其言，復改淮西行中書省為行樞密院}。

\noindent\textbf{元与南宋。}\eventanchor{SYD-1272-YUAN-073}{:爰自太祖皇帝以來，與宋使介交通}。

\subsection{1273年：元与安南、元、元与南宋}
《元史》、《安南志略》为这一年留下了可回查的年序入口。

\noindent\textbf{元。}\eventanchor{SYD-1273-YUAN-074}{帝以農事有益，詔勿禁}；\eventanchor{SYD-1273-YUAN-218}{庚寅，河南水，發粟賑民饑，仍免今年田租}；\eventanchor{SYD-1273-YUAN-290}{乙亥，詔：「免民代輸簽軍戶絲銀，及伐木夫戶賦稅}。

\noindent\textbf{元与安南。}\eventanchor{SYD-1273-YUAN-002}{安南使者還，言陳光昞受詔不拜}。

\noindent\textbf{元与南宋。}\eventanchor{SYD-1273-YUAN-146}{帝曰：「汝言雖是，若坐視宋人戍之，罪亦不免也}。

\subsection{1274年：元、元与高丽}
《元史》、《高丽史》为这一年留下了可回查的年序入口。

\noindent\textbf{元。}\eventanchor{SYD-1274-YUAN-003}{丁亥，罷各路洞冶總管府，以轉運司兼領}；\eventanchor{SYD-1274-YUAN-075}{改御史臺典事為都事}；\eventanchor{SYD-1274-YUAN-147}{命陝西等路宣撫使趙良弼為祕書監，充國信使，使日本}；\eventanchor{SYD-1274-YUAN-219}{世子愖奏乞隨朝及尚主，不許，命隨其父還國}。

\noindent\textbf{元与高丽。}\eventanchor{SYD-1274-YUAN-291}{又詔令國王頭輦哥等舉軍入高麗舊京，以脫朵兒、焦天翼為其國達魯花赤，護送禃還國，仍下詔：「林衍廢立，罪不可赦}。

\subsection{1275年：元、元与高丽、元与南宋}
《元史》、《高丽史》为这一年留下了可回查的年序入口。

\noindent\textbf{元。}\eventanchor{SYD-1275-YUAN-004}{:誕膺景命，奄四海以宅尊}；\eventanchor{SYD-1275-YUAN-076}{:我太祖聖武皇帝，握乾符而起朔土，以神武而膺帝圖，四震天聲，大恢土宇，輿圖之廣，歷古所無}；\eventanchor{SYD-1275-YUAN-148}{丙子，改山東、河間、陝西三路鹽課都轉運司為都轉運鹽使司}。

\noindent\textbf{元与高丽。}\eventanchor{SYD-1275-YUAN-220}{高麗國王王禃遣使奉表，為世子愖請昏}。

\noindent\textbf{元与南宋。}\eventanchor{SYD-1275-YUAN-292}{己亥，詔招諭宋襄陽守臣呂文煥}。

\subsection{1276年：元、元与高丽}
《元史》、《高丽史》为这一年留下了可回查的年序入口。

\noindent\textbf{元。}\eventanchor{SYD-1276-YUAN-005}{敕令整為書復之，賞整，使還軍中，誅永寧僧及其黨友}；\eventanchor{SYD-1276-YUAN-077}{敕燕王遣使持香旛，祠岳瀆、后土、五臺興國寺}；\eventanchor{SYD-1276-YUAN-149}{癸酉，同僉河南省事崔斌訟右丞阿里妄奏軍數二萬，敕杖而罷之}；\eventanchor{SYD-1276-YUAN-293}{詔自今凡詔令並以蒙古字行，仍遣百官子弟入學}。

\noindent\textbf{元与高丽。}\eventanchor{SYD-1276-YUAN-221}{壬辰，高麗國王王禃遣其臣齊安侯王淑來賀改國號}。

\subsection{1278年：元}
《元史》为这一年留下了可回查的年序入口。

\noindent\textbf{元。}\eventanchor{SYD-1278-YUAN-006}{近侍劉鐵木兒因言：「阿里海牙屬吏張鼎，今亦參知攻事}；\eventanchor{SYD-1278-YUAN-078}{申嚴無籍軍虜掠及傭奴代軍之禁}；\eventanchor{SYD-1278-YUAN-150}{戊申，從阿合馬請，自今御史臺非白于省，毋擅召倉庫吏，亦毋究錢穀數，及集議中書不至者罪之}；\eventanchor{SYD-1278-YUAN-222}{乙亥，敕省、院、臺諸司應聞奏事，必由起居注}；\eventanchor{SYD-1278-YUAN-294}{詔以也速海牙總制之}。

\subsection{1279年：元}
《元史》为这一年留下了可回查的年序入口。

\noindent\textbf{元。}\eventanchor{SYD-1279-YUAN-007}{西川軍官父死子繼，勤勞四十年，乞顯加爵秩}；\eventanchor{SYD-1279-YUAN-079}{甲戌，給要束合所領工匠牛二千，就令運米二千石供軍}；\eventanchor{SYD-1279-YUAN-151}{明日，下詔皇太子燕王參決朝政，凡中書省、樞密院、御史臺及百司之事，皆先啟後聞}；\eventanchor{SYD-1279-YUAN-223}{女直、水達達軍不出征者，令隸民籍輸賦}；\eventanchor{SYD-1279-YUAN-295}{數斬唐兀帶、武伯不花，餘減死論，以所掠者還其民}。

\subsection{1280年：元与高丽、安南及地方政权（年序桥接）、元}
《元史》、《高丽史》、《安南志略》为这一年留下了可回查的年序入口。

\noindent\textbf{元。}\eventanchor{SYD-1280-YUAN-008}{甲戌，給要束合所領工匠牛二千，就令運米二千石供軍}；\eventanchor{SYD-1280-YUAN-080}{近侍劉鐵木兒因言：「阿里海牙屬吏張鼎，今亦參知攻事}；\eventanchor{SYD-1280-YUAN-152}{明日，下詔皇太子燕王參決朝政，凡中書省、樞密院、御史臺及百司之事，皆先啟後聞}；\eventanchor{SYD-1280-YUAN-224}{女直、水達達軍不出征者，令隸民籍輸賦}；\eventanchor{SYD-1280-YUAN-296}{申嚴無籍軍虜掠及傭奴代軍之禁}。

\noindent\textbf{元与高丽、安南及地方政权（年序桥接）。}\eventanchor{SY-1280-CHRONOLOGY-BRIDGE}{【C级年序桥接】保留1280 年的多政权并行检索入口；本条不主张所列政权在当年发生直接互动，具体事件须另以单项材料入账}。

\subsection{1282年：元与高丽、安南及地方政权（年序桥接）}
《元史》、《高丽史》、《安南志略》为这一年留下了可回查的年序入口。

\noindent\textbf{元与高丽、安南及地方政权（年序桥接）。}\eventanchor{SY-1282-CHRONOLOGY-BRIDGE}{【C级年序桥接】保留1282 年的多政权并行检索入口；本条不主张所列政权在当年发生直接互动，具体事件须另以单项材料入账}。

\subsection{1283年：元与高丽、安南及地方政权（年序桥接）}
《元史》、《高丽史》、《安南志略》为这一年留下了可回查的年序入口。

\noindent\textbf{元与高丽、安南及地方政权（年序桥接）。}\eventanchor{SY-1283-CHRONOLOGY-BRIDGE}{【C级年序桥接】保留1283 年的多政权并行检索入口；本条不主张所列政权在当年发生直接互动，具体事件须另以单项材料入账}。

\subsection{1284年：元与高丽、安南及地方政权（年序桥接）}
《元史》、《高丽史》、《安南志略》为这一年留下了可回查的年序入口。

\noindent\textbf{元与高丽、安南及地方政权（年序桥接）。}\eventanchor{SY-1284-CHRONOLOGY-BRIDGE}{【C级年序桥接】保留1284 年的多政权并行检索入口；本条不主张所列政权在当年发生直接互动，具体事件须另以单项材料入账}。

\subsection{1285年：元与高丽、安南及地方政权（年序桥接）}
《元史》、《高丽史》、《安南志略》为这一年留下了可回查的年序入口。

\noindent\textbf{元与高丽、安南及地方政权（年序桥接）。}\eventanchor{SY-1285-CHRONOLOGY-BRIDGE}{【C级年序桥接】保留1285 年的多政权并行检索入口；本条不主张所列政权在当年发生直接互动，具体事件须另以单项材料入账}。

\subsection{1286年：元与高丽、安南及地方政权（年序桥接）}
《元史》、《高丽史》、《安南志略》为这一年留下了可回查的年序入口。

\noindent\textbf{元与高丽、安南及地方政权（年序桥接）。}\eventanchor{SY-1286-CHRONOLOGY-BRIDGE}{【C级年序桥接】保留1286 年的多政权并行检索入口；本条不主张所列政权在当年发生直接互动，具体事件须另以单项材料入账}。

\subsection{1288年：元与高丽、安南及地方政权（年序桥接）、元}
《元史》、《高丽史》、《安南志略》为这一年留下了可回查的年序入口。

\noindent\textbf{元。}\eventanchor{SYD-1288-YUAN-009}{:間者，行中書省右丞相伯顏遣使來奏，宋母后、幼主暨諸大臣百官，已於正月十八日齎璽綬奉表降附}；\eventanchor{SYD-1288-YUAN-081}{伯顏既受降表、玉璽，復遣囊加帶以趙尹甫、賈餘慶等還臨安，召宰相出議降事}；\eventanchor{SYD-1288-YUAN-153}{伯顏言：「張惠守宋府庫，不俟命擅啟管鑰}；\eventanchor{SYD-1288-YUAN-225}{帝命董文忠答之曰：「借使似道實輕汝曹，特似道一人之過耳，且汝主何負焉}；\eventanchor{SYD-1288-YUAN-297}{命以安慶、蘄、黃等郡宿兵，付宋都帶將之}。

\noindent\textbf{元与高丽、安南及地方政权（年序桥接）。}\eventanchor{SY-1288-CHRONOLOGY-BRIDGE}{【C级年序桥接】保留1288 年的多政权并行检索入口；本条不主张所列政权在当年发生直接互动，具体事件须另以单项材料入账}。

\subsection{1289年：元与高丽、安南及地方政权（年序桥接）、元}
《元史》、《高丽史》、《安南志略》为这一年留下了可回查的年序入口。

\noindent\textbf{元。}\eventanchor{SYD-1289-YUAN-010}{帝然其言，遽命止之}；\eventanchor{SYD-1289-YUAN-082}{壬子，寶應軍人施福殺其守將，降于淮東都元帥府，詔以福為千戶，佩金符}；\eventanchor{SYD-1289-YUAN-154}{詔賜鈔百八十錠贖還之}；\eventanchor{SYD-1289-YUAN-226}{詔鄭鼎所將侍衞軍萬人還京師，崔斌、阿里海牙同駐靜江，忽都鐵木兒、鄭鼎同駐鄂漢，賈居貞、脫博忽魯禿花同駐潭州}；\eventanchor{SYD-1289-YUAN-298}{昂吉兒、忻都、唐兀帶等引兵攻司空山寨，破之，殺張德興，執其三子以歸}。

\noindent\textbf{元与高丽、安南及地方政权（年序桥接）。}\eventanchor{SY-1289-CHRONOLOGY-BRIDGE}{【C级年序桥接】保留1289 年的多政权并行检索入口；本条不主张所列政权在当年发生直接互动，具体事件须另以单项材料入账}。

\subsection{1290年：元与高丽、安南及地方政权（年序桥接）}
《元史》、《高丽史》、《安南志略》为这一年留下了可回查的年序入口。

\noindent\textbf{元与高丽、安南及地方政权（年序桥接）。}\eventanchor{SY-1290-CHRONOLOGY-BRIDGE}{【C级年序桥接】保留1290 年的多政权并行检索入口；本条不主张所列政权在当年发生直接互动，具体事件须另以单项材料入账}。

\subsection{1291年：元与高丽、安南及地方政权（年序桥接）}
《元史》、《高丽史》、《安南志略》为这一年留下了可回查的年序入口。

\noindent\textbf{元与高丽、安南及地方政权（年序桥接）。}\eventanchor{SY-1291-CHRONOLOGY-BRIDGE}{【C级年序桥接】保留1291 年的多政权并行检索入口；本条不主张所列政权在当年发生直接互动，具体事件须另以单项材料入账}。

\subsection{1292年：元与高丽、安南及地方政权（年序桥接）}
《元史》、《高丽史》、《安南志略》为这一年留下了可回查的年序入口。

\noindent\textbf{元与高丽、安南及地方政权（年序桥接）。}\eventanchor{SY-1292-CHRONOLOGY-BRIDGE}{【C级年序桥接】保留1292 年的多政权并行检索入口；本条不主张所列政权在当年发生直接互动，具体事件须另以单项材料入账}。

\subsection{1293年：元与高丽、安南及地方政权（年序桥接）}
《元史》、《高丽史》、《安南志略》为这一年留下了可回查的年序入口。

\noindent\textbf{元与高丽、安南及地方政权（年序桥接）。}\eventanchor{SY-1293-CHRONOLOGY-BRIDGE}{【C级年序桥接】保留1293 年的多政权并行检索入口；本条不主张所列政权在当年发生直接互动，具体事件须另以单项材料入账}。

\subsection{1294年：元}
《元史》为这一年留下了可回查的年序入口。

\noindent\textbf{元。}\eventanchor{SYD-1294-YUAN-011}{:尚念先朝庶政，悉有成規，惟慎奉行，罔敢失墜}；\eventanchor{SYD-1294-YUAN-083}{:朕惟太祖聖武皇帝受天明命，肇造區夏，聖聖相承，光熙前緒}；\eventanchor{SYD-1294-YUAN-155}{帝命以軍職尊崇者授之}；\eventanchor{SYD-1294-YUAN-227}{定西平王奧魯赤、寧遠王闊闊出、鎮南王脫歡及也先帖木而大會賞賜例，金各五百兩、銀五千兩、鈔二千錠、幣帛各二百匹}；\eventanchor{SYD-1294-YUAN-299}{己未，復立平陽路之蒲、武鄉，保定路之博野，泰安州之新泰等縣}。

\subsection{1295年：元与高丽、安南及地方政权（年序桥接）、元}
《元史》、《高丽史》、《安南志略》为这一年留下了可回查的年序入口。

\noindent\textbf{元。}\eventanchor{SYD-1295-YUAN-012}{給江南行御史臺守護軍百人}；\eventanchor{SYD-1295-YUAN-084}{詔自今已後，專令中書擬奏}；\eventanchor{SYD-1295-YUAN-156}{〔己〕卯，罷四川淘金戶四千，還其元籍，罪初獻言者}；\eventanchor{SYD-1295-YUAN-228}{〔己〕卯，詔申飭中外：「有儒吏兼通者，各路舉之，廉訪司每道歲貢二人，省臺委官立法考試，中程者用之，所貢不公，罪其舉者}；\eventanchor{SYD-1295-YUAN-300}{罷河西軍，聽各還其所屬}。

\noindent\textbf{元与高丽、安南及地方政权（年序桥接）。}\eventanchor{SY-1295-CHRONOLOGY-BRIDGE}{【C级年序桥接】保留1295 年的多政权并行检索入口；本条不主张所列政权在当年发生直接互动，具体事件须另以单项材料入账}。

\subsection{1296年：元与高丽、安南及地方政权（年序桥接）、元、元与安南}
《元史》、《高丽史》、《安南志略》为这一年留下了可回查的年序入口。

\noindent\textbf{元。}\eventanchor{SYD-1296-YUAN-013}{甲戌，詔民間馬牛羊，百取其一，羊不滿百者亦取之，惟色目人及數乃取}；\eventanchor{SYD-1296-YUAN-085}{辛未，徙江浙行省拔都軍萬人戍潭州，潭州以南軍移戍郴州}；\eventanchor{SYD-1296-YUAN-157}{詔復立蒙樣剛等甸軍民官}；\eventanchor{SYD-1296-YUAN-229}{詔行省圖地形、覈軍實以聞}。

\noindent\textbf{元与安南。}\eventanchor{SYD-1296-YUAN-301}{安南國遣人招誘叛賊黃勝許}。

\noindent\textbf{元与高丽、安南及地方政权（年序桥接）。}\eventanchor{SY-1296-CHRONOLOGY-BRIDGE}{【C级年序桥接】保留1296 年的多政权并行检索入口；本条不主张所列政权在当年发生直接互动，具体事件须另以单项材料入账}。

\subsection{1297年：元与高丽、安南及地方政权（年序桥接）、元}
《元史》、《高丽史》、《安南志略》为这一年留下了可回查的年序入口。

\noindent\textbf{元。}\eventanchor{SYD-1297-YUAN-014}{帝曰：「礬課遣人覈實}；\eventanchor{SYD-1297-YUAN-086}{禁諸王、駙馬并權豪，毋奪民田，其獻田者有刑}；\eventanchor{SYD-1297-YUAN-158}{又賜緬王弟撒邦巴一珠虎符，酋領阿散三珠虎符，從者金符及金幣，遣之}；\eventanchor{SYD-1297-YUAN-230}{諸王也只里部忽剌帶於濟南商河縣侵擾居民，蹂踐禾稼，帝命詰之，走歸其部}；\eventanchor{SYD-1297-YUAN-302}{暗都剌火者所部四萬餘錠}。

\noindent\textbf{元与高丽、安南及地方政权（年序桥接）。}\eventanchor{SY-1297-CHRONOLOGY-BRIDGE}{【C级年序桥接】保留1297 年的多政权并行检索入口；本条不主张所列政权在当年发生直接互动，具体事件须另以单项材料入账}。

\subsection{1298年：元与高丽、安南及地方政权（年序桥接）、元}
《元史》、《高丽史》、《安南志略》为这一年留下了可回查的年序入口。

\noindent\textbf{元。}\eventanchor{SYD-1298-YUAN-015}{右丞相完澤言：「歲入之數，金一萬九千兩，銀六萬兩，鈔三百六十萬錠，然猶不足於用，又於至元鈔本中借二十萬錠}；\eventanchor{SYD-1298-YUAN-087}{〔衞〕輝、順德旱，大風損麥，免其田租一年}；\eventanchor{SYD-1298-YUAN-159}{罷建康金銀銅冶轉運司，還淘金戶於元籍，歲辦金悉責有司}；\eventanchor{SYD-1298-YUAN-231}{罷徐、邳爐冶所進息錢}；\eventanchor{SYD-1298-YUAN-303}{罷雲南行御史臺，置肅政廉訪司}。

\noindent\textbf{元与高丽、安南及地方政权（年序桥接）。}\eventanchor{SY-1298-CHRONOLOGY-BRIDGE}{【C级年序桥接】保留1298 年的多政权并行检索入口；本条不主张所列政权在当年发生直接互动，具体事件须另以单项材料入账}。

\subsection{1299年：元与高丽、安南及地方政权（年序桥接）、元}
《元史》、《高丽史》、《安南志略》为这一年留下了可回查的年序入口。

\noindent\textbf{元。}\eventanchor{SYD-1299-YUAN-016}{罷四川、福建等處行中書省，陝西行御史臺，江東、荊南、淮西三道宣慰司}；\eventanchor{SYD-1299-YUAN-088}{庚寅，詔遣使問民疾苦}；\eventanchor{SYD-1299-YUAN-160}{甲午，命何榮祖等更定律令}；\eventanchor{SYD-1299-YUAN-232}{丙申，揚州、淮安屬縣蝗，在地者為鶖啄食，飛者以翅擊死，詔禁捕鶖}；\eventanchor{SYD-1299-YUAN-304}{丙寅，詔各省戍軍輪次放還二年供役}。

\noindent\textbf{元与高丽、安南及地方政权（年序桥接）。}\eventanchor{SY-1299-CHRONOLOGY-BRIDGE}{【C级年序桥接】保留1299 年的多政权并行检索入口；本条不主张所列政权在当年发生直接互动，具体事件须另以单项材料入账}。

\subsection{1300年：元与高丽、安南及地方政权（年序桥接）、元}
《元史》、《高丽史》、《安南志略》为这一年留下了可回查的年序入口。

\noindent\textbf{元。}\eventanchor{SYD-1300-YUAN-017}{從之，仍賜交趾使人職名及弓矢鞍勒}；\eventanchor{SYD-1300-YUAN-089}{癸巳，詔諭和州諸城招集流移之民}；\eventanchor{SYD-1300-YUAN-161}{命和禮霍孫將兵與高和尚同赴北邊}；\eventanchor{SYD-1300-YUAN-233}{戊申，中書省臣議：「流通鈔法，凡賞賜宜多給幣帛，課程宜多收鈔}；\eventanchor{SYD-1300-YUAN-305}{詔以忙古帶仍行省福州}。

\noindent\textbf{元与高丽、安南及地方政权（年序桥接）。}\eventanchor{SY-1300-CHRONOLOGY-BRIDGE}{【C级年序桥接】保留1300 年的多政权并行检索入口；本条不主张所列政权在当年发生直接互动，具体事件须另以单项材料入账}。

\subsection{1301年：元与高丽、安南及地方政权（年序桥接）、元}
《元史》、《高丽史》、《安南志略》为这一年留下了可回查的年序入口。

\noindent\textbf{元。}\eventanchor{SYD-1301-YUAN-018}{癸未，命中書省會計姚演所領漣、海屯田官給之資與歲入之數，便則行之，否則罷去}；\eventanchor{SYD-1301-YUAN-090}{詔曰：「軍事卿等當自權衡之}；\eventanchor{SYD-1301-YUAN-162}{===校勘記===Category:元日戰爭‎}；\eventanchor{SYD-1301-YUAN-234}{阿塔海乞以戍三海口軍擊福建賊陳吊眼，詔以重勞不從}；\eventanchor{SYD-1301-YUAN-306}{八月甲子朔，招討使方文言擇守令、崇祀典、戢姦吏、禁盜賊、治軍旅、奬忠義六事，詔廷臣及諸老議舉行之}。

\noindent\textbf{元与高丽、安南及地方政权（年序桥接）。}\eventanchor{SY-1301-CHRONOLOGY-BRIDGE}{【C级年序桥接】保留1301 年的多政权并行检索入口；本条不主张所列政权在当年发生直接互动，具体事件须另以单项材料入账}。

\subsection{1302年：元与高丽、安南及地方政权（年序桥接）、元、元与安南}
《元史》、《高丽史》、《安南志略》为这一年留下了可回查的年序入口。

\noindent\textbf{元。}\eventanchor{SYD-1302-YUAN-019}{復請禁諸路釀酒，減免差稅，賑濟饑民}；\eventanchor{SYD-1302-YUAN-091}{命即行之，毋越三日}；\eventanchor{SYD-1302-YUAN-163}{〔己〕未，以諸王真童誣告濟南王，謫置劉國傑軍中自效}；\eventanchor{SYD-1302-YUAN-307}{八月甲子，詔御史臺凡有司婚姻、土田文案，遇赦依例檢覆}。

\noindent\textbf{元与安南。}\eventanchor{SYD-1302-YUAN-235}{安南國以馴象二及朱砂來獻}。

\noindent\textbf{元与高丽、安南及地方政权（年序桥接）。}\eventanchor{SY-1302-CHRONOLOGY-BRIDGE}{【C级年序桥接】保留1302 年的多政权并行检索入口；本条不主张所列政权在当年发生直接互动，具体事件须另以单项材料入账}。

\subsection{1303年：元与高丽、安南及地方政权（年序桥接）、元}
《元史》、《高丽史》、《安南志略》为这一年留下了可回查的年序入口。

\noindent\textbf{元。}\eventanchor{SYD-1303-YUAN-020}{罷江南財賦總管府及提舉司}；\eventanchor{SYD-1303-YUAN-092}{敕賜鈔五百錠，其稅課依例輸官}；\eventanchor{SYD-1303-YUAN-164}{戊子，弛太原、平陽酒禁}；\eventanchor{SYD-1303-YUAN-236}{中書省臣言：「宜捕送其所，令省、臺、宣政院遣官雜治}；\eventanchor{SYD-1303-YUAN-308}{安西轉運司於常課外增算五萬七千四百錠，人賜衣一襲，以勸其功}。

\noindent\textbf{元与高丽、安南及地方政权（年序桥接）。}\eventanchor{SY-1303-CHRONOLOGY-BRIDGE}{【C级年序桥接】保留1303 年的多政权并行检索入口；本条不主张所列政权在当年发生直接互动，具体事件须另以单项材料入账}。

\subsection{1304年：元与高丽、安南及地方政权（年序桥接）、元}
《元史》、《高丽史》、《安南志略》为这一年留下了可回查的年序入口。

\noindent\textbf{元。}\eventanchor{SYD-1304-YUAN-021}{敕軍民逃奴有獲者即付其主，主在他所者，赴所在官司給之，仍追逃奴鈔充獲者賞}；\eventanchor{SYD-1304-YUAN-093}{庚子，以永平、清、滄、柳林屯田被水，其逋租及民貸食者皆勿徵}；\eventanchor{SYD-1304-YUAN-165}{以平陽、太原去歲地大震，免其稅課一年}；\eventanchor{SYD-1304-YUAN-237}{自滎澤至睢州，築河防十有八所，給其夫鈔人十貫}；\eventanchor{SYD-1304-YUAN-309}{八年春正月己未，以災異故，詔天下恤民隱，省刑罰}。

\noindent\textbf{元与高丽、安南及地方政权（年序桥接）。}\eventanchor{SY-1304-CHRONOLOGY-BRIDGE}{【C级年序桥接】保留1304 年的多政权并行检索入口；本条不主张所列政权在当年发生直接互动，具体事件须另以单项材料入账}。

\subsection{1305年：元与高丽、安南及地方政权（年序桥接）、元}
《元史》、《高丽史》、《安南志略》为这一年留下了可回查的年序入口。

\noindent\textbf{元。}\eventanchor{SYD-1305-YUAN-022}{戊戌，詔芍陂、洪澤等屯田為豪右占據者，悉令輸租}；\eventanchor{SYD-1305-YUAN-094}{八月乙亥朔，省孛可孫冗員}；\eventanchor{SYD-1305-YUAN-166}{罷福建蒙古字提舉司及醫學提舉司}；\eventanchor{SYD-1305-YUAN-238}{寶慶路饑，發粟五千石賑之}；\eventanchor{SYD-1305-YUAN-310}{孛可孫專治芻粟，初惟數人，後以各位增入，遂至繁宂}。

\noindent\textbf{元与高丽、安南及地方政权（年序桥接）。}\eventanchor{SY-1305-CHRONOLOGY-BRIDGE}{【C级年序桥接】保留1305 年的多政权并行检索入口；本条不主张所列政权在当年发生直接互动，具体事件须另以单项材料入账}。

\subsection{1306年：元与高丽、安南及地方政权（年序桥接）、元}
《元史》、《高丽史》、《安南志略》为这一年留下了可回查的年序入口。

\noindent\textbf{元。}\eventanchor{SYD-1306-YUAN-023}{繼有旨，給瀛國公衣糧發遣之，唯與〔𤍟〕勿行}；\eventanchor{SYD-1306-YUAN-095}{有旨降民還之有司，征討所得，籍其數量，賜臣下有功者}；\eventanchor{SYD-1306-YUAN-167}{阿里海牙復鎮遠軍，發軍千人戍守，以其地與西川行省接，就以隸焉}；\eventanchor{SYD-1306-YUAN-239}{八月丁亥朔，給乾山造船軍匠冬衣，及新附軍鈔}；\eventanchor{SYD-1306-YUAN-311}{罷都功德使脫烈，其修設佛事妄費官物，皆徵還之}。

\noindent\textbf{元与高丽、安南及地方政权（年序桥接）。}\eventanchor{SY-1306-CHRONOLOGY-BRIDGE}{【C级年序桥接】保留1306 年的多政权并行检索入口；本条不主张所列政权在当年发生直接互动，具体事件须另以单项材料入账}。

\subsection{1307年：元}
《元史》为这一年留下了可回查的年序入口。

\noindent\textbf{元。}\eventanchor{SYD-1307-YUAN-024}{併省江淮、雲南州郡}；\eventanchor{SYD-1307-YUAN-096}{己未，吏部尚書劉好禮以吉利吉思風俗事宜來上}；\eventanchor{SYD-1307-YUAN-168}{甲寅，降太醫院為尚醫監，改給銅印}；\eventanchor{SYD-1307-YUAN-240}{前後衞軍自願征日本者，命選留五衞漢軍千餘，其新附軍令悉行}；\eventanchor{SYD-1307-YUAN-312}{時各衞議先遣七人，而以三人自代，從之}。

\subsection{1308年：元与高丽、安南及地方政权（年序桥接）、元}
《元史》、《高丽史》、《安南志略》为这一年留下了可回查的年序入口。

\noindent\textbf{元。}\eventanchor{SYD-1308-YUAN-025}{帝曰：「礬課遣人覈實}；\eventanchor{SYD-1308-YUAN-097}{甲戌，詔民間馬牛羊，百取其一，羊不滿百者亦取之，惟色目人及數乃取}；\eventanchor{SYD-1308-YUAN-169}{禁諸王、駙馬并權豪，毋奪民田，其獻田者有刑}；\eventanchor{SYD-1308-YUAN-241}{辛未，徙江浙行省拔都軍萬人戍潭州，潭州以南軍移戍郴州}；\eventanchor{SYD-1308-YUAN-313}{又賜緬王弟撒邦巴一珠虎符，酋領阿散三珠虎符，從者金符及金幣，遣之}。

\noindent\textbf{元与高丽、安南及地方政权（年序桥接）。}\eventanchor{SY-1308-CHRONOLOGY-BRIDGE}{【C级年序桥接】保留1308 年的多政权并行检索入口；本条不主张所列政权在当年发生直接互动，具体事件须另以单项材料入账}。

\subsection{1309年：元与高丽、安南及地方政权（年序桥接）、元}
《元史》、《高丽史》、《安南志略》为这一年留下了可回查的年序入口。

\noindent\textbf{元。}\eventanchor{SYD-1309-YUAN-026}{頒行至大銀鈔，詔曰：「昔我世祖皇帝既登大寶，始造中統交鈔，以便民用，歲久法隳，亦既更張，印造至元寶鈔}；\eventanchor{SYD-1309-YUAN-098}{大都立資國院，秩正二品}；\eventanchor{SYD-1309-YUAN-170}{帝曰：「大事也，其速擇使復齎璽書往招諭}；\eventanchor{SYD-1309-YUAN-242}{甲申，賜寧肅王脫脫金印}；\eventanchor{SYD-1309-YUAN-314}{壬午，江南行臺劾：「平章政事教化，詐言家貧，冒受賜貨物，折鈔二萬錠}。

\noindent\textbf{元与高丽、安南及地方政权（年序桥接）。}\eventanchor{SY-1309-CHRONOLOGY-BRIDGE}{【C级年序桥接】保留1309 年的多政权并行检索入口；本条不主张所列政权在当年发生直接互动，具体事件须另以单项材料入账}。

\subsection{1310年：元与高丽、安南及地方政权（年序桥接）、元}
《元史》、《高丽史》、《安南志略》为这一年留下了可回查的年序入口。

\noindent\textbf{元。}\eventanchor{SYD-1310-YUAN-027}{汴梁、懷孟、衞輝、彰德、歸德、汝寧、南陽、河南等路蝗}；\eventanchor{SYD-1310-YUAN-099}{敕臺臣遣官往鞫，毋徇私情}；\eventanchor{SYD-1310-YUAN-171}{從之，仍賜鈔三千錠賑其部貧民}；\eventanchor{SYD-1310-YUAN-243}{帝曰：「世祖謀慮深遠若是，待諸王朝會，頒賞既畢，卿等備述其故，然後與之，使彼知愧}；\eventanchor{SYD-1310-YUAN-315}{封僧亦憐真乞烈思為文國公，賜金印}。

\noindent\textbf{元与高丽、安南及地方政权（年序桥接）。}\eventanchor{SY-1310-CHRONOLOGY-BRIDGE}{【C级年序桥接】保留1310 年的多政权并行检索入口；本条不主张所列政权在当年发生直接互动，具体事件须另以单项材料入账}。

\subsection{1311年：元}
《元史》为这一年留下了可回查的年序入口。

\noindent\textbf{元。}\eventanchor{SYD-1311-YUAN-028}{壬午，靈駕發引，葬起輦谷，從諸帝陵}；\eventanchor{SYD-1311-YUAN-100}{四年春正月癸酉朔，帝不豫，免朝賀，大赦天下}；\eventanchor{SYD-1311-YUAN-172}{武宗當富有之大業，慨然欲創治改法而有為，故其封爵太盛，而遙授之官眾，錫賚太隆，而泛賞之恩溥，至元、大德之政，於是稍有變更云}；\eventanchor{SYD-1311-YUAN-244}{敕：「諸王、駙馬戶在縉山、懷來、永興縣者，與民均服徭役}；\eventanchor{SYD-1311-YUAN-316}{敕直省舍人以其半給事殿庭，半聽中書差遣}。

\subsection{1312年：元与高丽、安南及地方政权（年序桥接）、元}
《元史》、《高丽史》、《安南志略》为这一年留下了可回查的年序入口。

\noindent\textbf{元。}\eventanchor{SYD-1312-YUAN-029}{己卯，馬八兒國遣使貢珍珠、異寶、縑段}；\eventanchor{SYD-1312-YUAN-101}{己未，罷雲南都元帥府，所管軍民隸行省}；\eventanchor{SYD-1312-YUAN-173}{己未，荊湖占城行省言：「忽都虎、忽馬兒等將兵征占城，前鋒舟師至舒眉蓮港不知所向，令萬戶劉君慶進軍次新州，獲占蠻，始知我軍已還矣}；\eventanchor{SYD-1312-YUAN-245}{詔依所言汰選，毋徇私情}；\eventanchor{SYD-1312-YUAN-317}{詔以中書右丞、行省事忙兀台為江淮等處行中書省平章政事，其行省左丞忽剌出、蒲壽庚，參政管如德分省泉州}。

\noindent\textbf{元与高丽、安南及地方政权（年序桥接）。}\eventanchor{SY-1312-CHRONOLOGY-BRIDGE}{【C级年序桥接】保留1312 年的多政权并行检索入口；本条不主张所列政权在当年发生直接互动，具体事件须另以单项材料入账}。

\subsection{1313年：元}
《元史》为这一年留下了可回查的年序入口。

\noindent\textbf{元。}\eventanchor{SYD-1313-YUAN-030}{安童與諸老臣議世榮所行，當罷者罷之，更者更之，其所用人實無罪者，朕自裁決}；\eventanchor{SYD-1313-YUAN-102}{丙子，真蠟、占城貢樂工十人及藥材、鰐魚皮諸物}；\eventanchor{SYD-1313-YUAN-174}{癸亥，敕以麥朮丁所行清潔，與安童治省事}；\eventanchor{SYD-1313-YUAN-246}{江淮行省以戰船千艘習水戰江中}；\eventanchor{SYD-1313-YUAN-318}{盧世榮請罷福建行中書省，立宣慰司，隸江西行中書省}。

\subsection{1314年：元与高丽、安南及地方政权（年序桥接）、元}
《元史》、《高丽史》、《安南志略》为这一年留下了可回查的年序入口。

\noindent\textbf{元。}\eventanchor{SYD-1314-YUAN-031}{大寧路地震，有聲如雷}；\eventanchor{SYD-1314-YUAN-103}{帝曰：「如更不悛，則罷不敍}；\eventanchor{SYD-1314-YUAN-175}{衡州、郴州、興國、永州路、耒陽州饑，發廪減價賑糶}；\eventanchor{SYD-1314-YUAN-247}{遣官括淮民所佃閑田不輸稅者}；\eventanchor{SYD-1314-YUAN-319}{申飭內侍及諸司隔越中書奏請之禁}。

\noindent\textbf{元与高丽、安南及地方政权（年序桥接）。}\eventanchor{SY-1314-CHRONOLOGY-BRIDGE}{【C级年序桥接】保留1314 年的多政权并行检索入口；本条不主张所列政权在当年发生直接互动，具体事件须另以单项材料入账}。

\subsection{1315年：元与高丽、安南及地方政权（年序桥接）、元}
《元史》、《高丽史》、《安南志略》为这一年留下了可回查的年序入口。

\noindent\textbf{元。}\eventanchor{SYD-1315-YUAN-032}{帝曰：「此朕之愆，豈卿等所致，其復乃職}；\eventanchor{SYD-1315-YUAN-104}{詔明言其事當行者以聞}；\eventanchor{SYD-1315-YUAN-176}{八月丙戌，贛州賊蔡五九陷汀州寧花縣，僭稱王號，詔遣江浙行省平章張驢等率兵討之}；\eventanchor{SYD-1315-YUAN-248}{丙申，賜諸王納忽答兒金五十兩、銀二百五十兩、鈔五百錠}；\eventanchor{SYD-1315-YUAN-320}{丙午，封諸王察八兒為汝寧王}。

\noindent\textbf{元与高丽、安南及地方政权（年序桥接）。}\eventanchor{SY-1315-CHRONOLOGY-BRIDGE}{【C级年序桥接】保留1315 年的多政权并行检索入口；本条不主张所列政权在当年发生直接互动，具体事件须另以单项材料入账}。

\subsection{1316年：元与高丽、安南及地方政权（年序桥接）、元}
《元史》、《高丽史》、《安南志略》为这一年留下了可回查的年序入口。

\noindent\textbf{元。}\eventanchor{SYD-1316-YUAN-033}{罷江南財賦總管府及提舉司}；\eventanchor{SYD-1316-YUAN-105}{敕賜鈔五百錠，其稅課依例輸官}；\eventanchor{SYD-1316-YUAN-177}{敕軍民逃奴有獲者即付其主，主在他所者，赴所在官司給之，仍追逃奴鈔充獲者賞}；\eventanchor{SYD-1316-YUAN-249}{庚子，以永平、清、滄、柳林屯田被水，其逋租及民貸食者皆勿徵}；\eventanchor{SYD-1316-YUAN-321}{戊戌，詔芍陂、洪澤等屯田為豪右占據者，悉令輸租}。

\noindent\textbf{元与高丽、安南及地方政权（年序桥接）。}\eventanchor{SY-1316-CHRONOLOGY-BRIDGE}{【C级年序桥接】保留1316 年的多政权并行检索入口；本条不主张所列政权在当年发生直接互动，具体事件须另以单项材料入账}。

\subsection{1317年：元与高丽、安南及地方政权（年序桥接）、元}
《元史》、《高丽史》、《安南志略》为这一年留下了可回查的年序入口。

\noindent\textbf{元。}\eventanchor{SYD-1317-YUAN-034}{敕免雲南從征交趾蒙古軍屯田租}；\eventanchor{SYD-1317-YUAN-106}{帝曰：「比遣人往，事已緩矣}；\eventanchor{SYD-1317-YUAN-178}{帝曰：「苟事力未損，卽遣之}；\eventanchor{SYD-1317-YUAN-250}{帝曰：「汝漢人用事者，豈皆賢邪？」江南諸路學田昔皆隸官，詔復給本學，以便教養}；\eventanchor{SYD-1317-YUAN-322}{帝曰：「朕以漢人徇私，用泰和律處事，致盜賊滋眾，故有是言}。

\noindent\textbf{元与高丽、安南及地方政权（年序桥接）。}\eventanchor{SY-1317-CHRONOLOGY-BRIDGE}{【C级年序桥接】保留1317 年的多政权并行检索入口；本条不主张所列政权在当年发生直接互动，具体事件须另以单项材料入账}。

\subsection{1318年：元与高丽、安南及地方政权（年序桥接）、元}
《元史》、《高丽史》、《安南志略》为这一年留下了可回查的年序入口。

\noindent\textbf{元。}\eventanchor{SYD-1318-YUAN-035}{帝皆從其議，詔行之}；\eventanchor{SYD-1318-YUAN-107}{帝曰：「除陳巖、呂師夔、管如德、范文虎四人，餘從卿議}；\eventanchor{SYD-1318-YUAN-179}{帝曰：「囚非羣羊，豈可遽殺耶！宜悉配隸淘金}；\eventanchor{SYD-1318-YUAN-251}{帝自將征乃顏，發上都}；\eventanchor{SYD-1318-YUAN-323}{詔從之，各設官六員}。

\noindent\textbf{元与高丽、安南及地方政权（年序桥接）。}\eventanchor{SY-1318-CHRONOLOGY-BRIDGE}{【C级年序桥接】保留1318 年的多政权并行检索入口；本条不主张所列政权在当年发生直接互动，具体事件须另以单项材料入账}。

\subsection{1319年：元与高丽、安南及地方政权（年序桥接）、元}
《元史》、《高丽史》、《安南志略》为这一年留下了可回查的年序入口。

\noindent\textbf{元。}\eventanchor{SYD-1319-YUAN-036}{丙寅，改懷孟路為懷慶路}；\eventanchor{SYD-1319-YUAN-108}{敕：「諸司有受命不之官及避繁劇託故去職者，奪其宣敕}；\eventanchor{SYD-1319-YUAN-180}{敕上都、大都冬夏設食于路，以食饑者}；\eventanchor{SYD-1319-YUAN-252}{封諸王月魯鐵木兒為恩王，給印，置王傅官}；\eventanchor{SYD-1319-YUAN-324}{河間民饑，發粟賑之}。

\noindent\textbf{元与高丽、安南及地方政权（年序桥接）。}\eventanchor{SY-1319-CHRONOLOGY-BRIDGE}{【C级年序桥接】保留1319 年的多政权并行检索入口；本条不主张所列政权在当年发生直接互动，具体事件须另以单项材料入账}。

\subsection{1320年：元}
《元史》为这一年留下了可回查的年序入口。

\noindent\textbf{元。}\eventanchor{SYD-1320-YUAN-037}{癸未，帝御大明殿，受諸王、百官朝賀}；\eventanchor{SYD-1320-YUAN-109}{中書省臣進曰：「臺臣所言良是，若非振理朝綱，法度愈壞}；\eventanchor{SYD-1320-YUAN-181}{臣等乞賜罷黜，選任賢者}；\eventanchor{SYD-1320-YUAN-253}{壬午，御史臺臣言：「比賜不兒罕丁山場、完者不花海舶稅，會計其鈔，皆數十萬錠，諸王軍民貧乏者，所賜未嘗若是，苟不撙節，漸致帑藏虛竭，民益困矣}；\eventanchor{SYD-1320-YUAN-325}{太官進膳，必分賜貴近}。

\subsection{1321年：元与高丽、安南及地方政权（年序桥接）、元}
《元史》、《高丽史》、《安南志略》为这一年留下了可回查的年序入口。

\noindent\textbf{元。}\eventanchor{SYD-1321-YUAN-038}{成宗大德三年，以寧遠王闊闊出總兵北邊，怠於備禦，命帝即軍中代之}；\eventanchor{SYD-1321-YUAN-110}{海都不得志去，旋亦死}；\eventanchor{SYD-1321-YUAN-182}{明日復戰，軍少却，海都乘之，帝揮軍力戰，突出敵陣後，全軍而還}；\eventanchor{SYD-1321-YUAN-254}{師失利，親出陣，力戰大敗之，盡獲其輜重，悉援諸王、駙馬眾軍以出}；\eventanchor{SYD-1321-YUAN-326}{十二月，軍至按台山，乃蠻帶部落降}。

\noindent\textbf{元与高丽、安南及地方政权（年序桥接）。}\eventanchor{SY-1321-CHRONOLOGY-BRIDGE}{【C级年序桥接】保留1321 年的多政权并行检索入口；本条不主张所列政权在当年发生直接互动，具体事件须另以单项材料入账}。

\subsection{1322年：元与高丽、安南及地方政权（年序桥接）、元}
《元史》、《高丽史》、《安南志略》为这一年留下了可回查的年序入口。

\noindent\textbf{元。}\eventanchor{SYD-1322-YUAN-039}{帝曰：「父兄雖死事，子弟不勝任者，安可用之？苟賢矣，則病故者亦不可降也}；\eventanchor{SYD-1322-YUAN-111}{帝曰：「自今不當給者汝即畫之，當給者宜覆奏，朕自處之}；\eventanchor{SYD-1322-YUAN-183}{甲申，詔皇孫撫諸軍討叛王火魯火孫、合丹禿魯干}；\eventanchor{SYD-1322-YUAN-255}{募民能耕江南曠土及公田者，免其差役三年，其輸租免三分之一}；\eventanchor{SYD-1322-YUAN-327}{辛卯，以六衞漢兵千二百、新附軍四百、屯田兵四百造尚書省}。

\noindent\textbf{元与高丽、安南及地方政权（年序桥接）。}\eventanchor{SY-1322-CHRONOLOGY-BRIDGE}{【C级年序桥接】保留1322 年的多政权并行检索入口；本条不主张所列政权在当年发生直接互动，具体事件须另以单项材料入账}。

\subsection{1323年：元、元与安南}
《元史》、《安南志略》为这一年留下了可回查的年序入口。

\noindent\textbf{元。}\eventanchor{SYD-1323-YUAN-040}{帝曰：「汝家不與它漢人比，弓矢不汝禁也，任汝執之}；\eventanchor{SYD-1323-YUAN-184}{皇孫甘不剌所部軍乏食，發大同路榷場糧賑之}；\eventanchor{SYD-1323-YUAN-256}{命贛州路發米千八百九十石賑之}；\eventanchor{SYD-1323-YUAN-328}{桑哥以聞，請擢絜矩為尚書省舍人，從之}。

\noindent\textbf{元与安南。}\eventanchor{SYD-1323-YUAN-112}{丁亥，安南國王陳日烜遣使來貢方物}。

\subsection{1324年：元与高丽、安南及地方政权（年序桥接）、元}
《元史》、《高丽史》、《安南志略》为这一年留下了可回查的年序入口。

\noindent\textbf{元。}\eventanchor{SYD-1324-YUAN-041}{頒行至大銀鈔，詔曰：「昔我世祖皇帝既登大寶，始造中統交鈔，以便民用，歲久法隳，亦既更張，印造至元寶鈔}；\eventanchor{SYD-1324-YUAN-113}{汴梁、懷孟、衞輝、彰德、歸德、汝寧、南陽、河南等路蝗}；\eventanchor{SYD-1324-YUAN-185}{敕臺臣遣官往鞫，毋徇私情}；\eventanchor{SYD-1324-YUAN-257}{從之，仍賜鈔三千錠賑其部貧民}；\eventanchor{SYD-1324-YUAN-329}{大都立資國院，秩正二品}。

\noindent\textbf{元与高丽、安南及地方政权（年序桥接）。}\eventanchor{SY-1324-CHRONOLOGY-BRIDGE}{【C级年序桥接】保留1324 年的多政权并行检索入口；本条不主张所列政权在当年发生直接互动，具体事件须另以单项材料入账}。

\subsection{1325年：元与高丽、安南及地方政权（年序桥接）}
《元史》、《高丽史》、《安南志略》为这一年留下了可回查的年序入口。

\noindent\textbf{元与高丽、安南及地方政权（年序桥接）。}\eventanchor{SY-1325-CHRONOLOGY-BRIDGE}{【C级年序桥接】保留1325 年的多政权并行检索入口；本条不主张所列政权在当年发生直接互动，具体事件须另以单项材料入账}。

\subsection{1326年：元与高丽、安南及地方政权（年序桥接）}
《元史》、《高丽史》、《安南志略》为这一年留下了可回查的年序入口。

\noindent\textbf{元与高丽、安南及地方政权（年序桥接）。}\eventanchor{SY-1326-CHRONOLOGY-BRIDGE}{【C级年序桥接】保留1326 年的多政权并行检索入口；本条不主张所列政权在当年发生直接互动，具体事件须另以单项材料入账}。

\subsection{1327年：元与高丽、安南及地方政权（年序桥接）、元、元与安南}
《元史》、《高丽史》、《安南志略》为这一年留下了可回查的年序入口。

\noindent\textbf{元。}\eventanchor{SYD-1327-YUAN-042}{敕遼陽行省驗實給之}；\eventanchor{SYD-1327-YUAN-114}{帝曰：「此亦何待上聞，當速賑之！」凡出粟五十八萬二千八百八十九石}；\eventanchor{SYD-1327-YUAN-186}{平灤民萬五千四百六十五戶饑，賑粟五千石}；\eventanchor{SYD-1327-YUAN-258}{阿速敦等二百九十五人乏食，命驗其實，給糧賑之}。

\noindent\textbf{元与安南。}\eventanchor{SYD-1327-YUAN-330}{安南國王陳日烜遣其中大夫陳克用來貢方物}。

\noindent\textbf{元与高丽、安南及地方政权（年序桥接）。}\eventanchor{SY-1327-CHRONOLOGY-BRIDGE}{【C级年序桥接】保留1327 年的多政权并行检索入口；本条不主张所列政权在当年发生直接互动，具体事件须另以单项材料入账}。

\subsection{1328年：元}
《元史》为这一年留下了可回查的年序入口。

\noindent\textbf{元。}\eventanchor{SYD-1328-YUAN-043}{罷江南諸省買銀提舉司}；\eventanchor{SYD-1328-YUAN-115}{帝曰：「桑哥已誅，納〔速〕剌丁滅里在獄，唯沙不丁朕姑釋之耳}；\eventanchor{SYD-1328-YUAN-187}{何榮祖言：「宜召各省官任錢穀者詣京師，集議科取之法以聞}；\eventanchor{SYD-1328-YUAN-259}{壬申，立河南江北行中書省，治汴梁}；\eventanchor{SYD-1328-YUAN-331}{塔叉兒、塔帶民饑，發米賑之}。

\subsection{1329年：元与高丽、安南及地方政权（年序桥接）、元}
《元史》、《高丽史》、《安南志略》为这一年留下了可回查的年序入口。

\noindent\textbf{元。}\eventanchor{SYD-1329-YUAN-044}{丙寅，改懷孟路為懷慶路}；\eventanchor{SYD-1329-YUAN-116}{播州南寧長官洛麼作亂，思州守臣換住哥招諭之，洛麼遣人以方物來覲}；\eventanchor{SYD-1329-YUAN-188}{敕：「諸司有受命不之官及避繁劇託故去職者，奪其宣敕}；\eventanchor{SYD-1329-YUAN-260}{敕上都、大都冬夏設食于路，以食饑者}。

\noindent\textbf{元与高丽、安南及地方政权（年序桥接）。}\eventanchor{SY-1329-CHRONOLOGY-BRIDGE}{【C级年序桥接】保留1329 年的多政权并行检索入口；本条不主张所列政权在当年发生直接互动，具体事件须另以单项材料入账}。

\subsection{1330年：元与高丽、安南及地方政权（年序桥接）、元}
《元史》、《高丽史》、《安南志略》为这一年留下了可回查的年序入口。

\noindent\textbf{元。}\eventanchor{SYD-1330-YUAN-045}{帝謂中書省臣曰：「至尊委我以天下事，日夜寅畏，惟恐弗堪}；\eventanchor{SYD-1330-YUAN-117}{監察御史段輔、太子詹事郭貫等，首請近賢人，擇師傅，帝嘉納之}；\eventanchor{SYD-1330-YUAN-189}{六年十月戊午，受玉冊，詔命百司庶務必先啟太子，然後奏聞}；\eventanchor{SYD-1330-YUAN-261}{仁宗欲立為太子，帝入謁太后固辭，曰：「臣幼無能，且有兄在，宜立兄，以臣輔之}。

\noindent\textbf{元与高丽、安南及地方政权（年序桥接）。}\eventanchor{SY-1330-CHRONOLOGY-BRIDGE}{【C级年序桥接】保留1330 年的多政权并行检索入口；本条不主张所列政权在当年发生直接互动，具体事件须另以单项材料入账}。

\subsection{1331年：元与高丽、安南及地方政权（年序桥接）}
《元史》、《高丽史》、《安南志略》为这一年留下了可回查的年序入口。

\noindent\textbf{元与高丽、安南及地方政权（年序桥接）。}\eventanchor{SY-1331-CHRONOLOGY-BRIDGE}{【C级年序桥接】保留1331 年的多政权并行检索入口；本条不主张所列政权在当年发生直接互动，具体事件须另以单项材料入账}。

\subsection{1332年：元与高丽、安南及地方政权（年序桥接）、元}
《元史》、《高丽史》、《安南志略》为这一年留下了可回查的年序入口。

\noindent\textbf{元。}\eventanchor{SYD-1332-YUAN-046}{丙午，中書省臣言：「京畿荐饑，宜免今歲田租}；\eventanchor{SYD-1332-YUAN-118}{敕依東平例，發附近官廪，計口以給}；\eventanchor{SYD-1332-YUAN-190}{敕應昌府給乞答帶糧五百石，以賑饑民}；\eventanchor{SYD-1332-YUAN-262}{帝曰：「死者勿論，其存者罰不可恕也}。

\noindent\textbf{元与高丽、安南及地方政权（年序桥接）。}\eventanchor{SY-1332-CHRONOLOGY-BRIDGE}{【C级年序桥接】保留1332 年的多政权并行检索入口；本条不主张所列政权在当年发生直接互动，具体事件须另以单项材料入账}。

\subsection{1333年：元、元与高丽}
《元史》、《高丽史》为这一年留下了可回查的年序入口。

\noindent\textbf{元。}\eventanchor{SYD-1333-YUAN-047}{敕海運米十萬石給遼陽戍兵，仍諭其省官薛闍干，令伯鐵木部欽察等耕漁自養，糧不須給}；\eventanchor{SYD-1333-YUAN-119}{帝曰：「非其人不行，乃朕中止之耳，勿徵}；\eventanchor{SYD-1333-YUAN-263}{甲戌，河南江北行省平章伯顏言：「揚州忙兀臺所立屯田，為田四萬餘頃，官種外，宜聽民耕墾}。

\noindent\textbf{元与高丽。}\eventanchor{SYD-1333-YUAN-191}{高麗國王王〔賰〕請易名曰昛，其僉議府請陞僉議司，降二品印，從之}。

\subsection{1334年：元与高丽、安南及地方政权（年序桥接）、元}
《元史》、《高丽史》、《安南志略》为这一年留下了可回查的年序入口。

\noindent\textbf{元。}\eventanchor{SYD-1334-YUAN-048}{甲辰，遣司徒兀都帶、平章政事不忽木、左丞張九思，率百官請諡于南郊}；\eventanchor{SYD-1334-YUAN-120}{親王、諸大臣發使告哀于皇孫}；\eventanchor{SYD-1334-YUAN-192}{三十一年春正月壬子朔，帝不豫，免朝賀}；\eventanchor{SYD-1334-YUAN-264}{世祖度量弘廣，知人善任使，信用儒術，用能以夏變夷，立經陳紀，所以為一代之制者，規模宏遠矣}。

\noindent\textbf{元与高丽、安南及地方政权（年序桥接）。}\eventanchor{SY-1334-CHRONOLOGY-BRIDGE}{【C级年序桥接】保留1334 年的多政权并行检索入口；本条不主张所列政权在当年发生直接互动，具体事件须另以单项材料入账}。

\subsection{1335年：元与高丽、安南及地方政权（年序桥接）、元}
《元史》、《高丽史》、《安南志略》为这一年留下了可回查的年序入口。

\noindent\textbf{元。}\eventanchor{SYD-1335-YUAN-049}{丁丑，頒農桑舊制十四條于天下，仍詔勵有司以察勤惰}；\eventanchor{SYD-1335-YUAN-121}{甲申，廣西花角蠻為寇，命所部討之}；\eventanchor{SYD-1335-YUAN-193}{禁僧道買民田，違者坐罪，沒其直}；\eventanchor{SYD-1335-YUAN-265}{乙亥，賜公主不答昔你媵戶鈔四千錠}。

\noindent\textbf{元与高丽、安南及地方政权（年序桥接）。}\eventanchor{SY-1335-CHRONOLOGY-BRIDGE}{【C级年序桥接】保留1335 年的多政权并行检索入口；本条不主张所列政权在当年发生直接互动，具体事件须另以单项材料入账}。

\subsection{1336年：元与高丽、安南及地方政权（年序桥接）、元、元与安南}
《元史》、《高丽史》、《安南志略》为这一年留下了可回查的年序入口。

\noindent\textbf{元。}\eventanchor{SYD-1336-YUAN-050}{:洪惟太祖皇帝膺期撫運，肇開帝業}；\eventanchor{SYD-1336-YUAN-194}{帝問曰：「所賜為誰？」對曰：「左丞相阿散所得為多}；\eventanchor{SYD-1336-YUAN-266}{帝曰：「言事者直至朕前可也，如細民輒訴訟者則禁之}。

\noindent\textbf{元与安南。}\eventanchor{SYD-1336-YUAN-122}{安南國遣其臣鄧恭儉來貢方物}。

\noindent\textbf{元与高丽、安南及地方政权（年序桥接）。}\eventanchor{SY-1336-CHRONOLOGY-BRIDGE}{【C级年序桥接】保留1336 年的多政权并行检索入口；本条不主张所列政权在当年发生直接互动，具体事件须另以单项材料入账}。

\subsection{1337年：元与高丽、安南及地方政权（年序桥接）}
《元史》、《高丽史》、《安南志略》为这一年留下了可回查的年序入口。

\noindent\textbf{元与高丽、安南及地方政权（年序桥接）。}\eventanchor{SY-1337-CHRONOLOGY-BRIDGE}{【C级年序桥接】保留1337 年的多政权并行检索入口；本条不主张所列政权在当年发生直接互动，具体事件须另以单项材料入账}。

\subsection{1338年：元与高丽、安南及地方政权（年序桥接）、元、元与安南}
《元史》、《高丽史》、《安南志略》为这一年留下了可回查的年序入口。

\noindent\textbf{元。}\eventanchor{SYD-1338-YUAN-051}{敕：「諸王、駙馬戶在縉山、懷來、永興縣者，與民均服徭役}；\eventanchor{SYD-1338-YUAN-123}{敕直省舍人以其半給事殿庭，半聽中書差遣}；\eventanchor{SYD-1338-YUAN-195}{帝納其言，凡營繕悉罷之}。

\noindent\textbf{元与安南。}\eventanchor{SYD-1338-YUAN-267}{帝謂省臣曰：「安南國王慕義來歸，宜厚其賜，以懷遠人，其進勳爵、受田如故}。

\noindent\textbf{元与高丽、安南及地方政权（年序桥接）。}\eventanchor{SY-1338-CHRONOLOGY-BRIDGE}{【C级年序桥接】保留1338 年的多政权并行检索入口；本条不主张所列政权在当年发生直接互动，具体事件须另以单项材料入账}。

\subsection{1339年：元与高丽、安南及地方政权（年序桥接）、元}
《元史》、《高丽史》、《安南志略》为这一年留下了可回查的年序入口。

\noindent\textbf{元。}\eventanchor{SYD-1339-YUAN-052}{:大都、上都、興和三路，差稅免三年}；\eventanchor{SYD-1339-YUAN-124}{丙寅，楚丘縣河堤壞，發民丁二千三百五十人修之}；\eventanchor{SYD-1339-YUAN-196}{賓天之日，皇后傳顧命於太師、太平王、右丞相、答剌罕燕帖木兒，太保、浚寧王、知樞密院事伯顏等，謂聖體彌留，益推固讓之初志，以宗社之重，屬諸大兄忽都篤皇帝之世嫡}；\eventanchor{SYD-1339-YUAN-268}{丙辰，給宿衞士、蒙古、漢軍三萬人禦寒衣}。

\noindent\textbf{元与高丽、安南及地方政权（年序桥接）。}\eventanchor{SY-1339-CHRONOLOGY-BRIDGE}{【C级年序桥接】保留1339 年的多政权并行检索入口；本条不主张所列政权在当年发生直接互动，具体事件须另以单项材料入账}。

\subsection{1340年：元}
《元史》为这一年留下了可回查的年序入口。

\noindent\textbf{元。}\eventanchor{SYD-1340-YUAN-053}{敕更賜璽書三十二、圓符四，仍究詰所失者}；\eventanchor{SYD-1340-YUAN-125}{庚辰，以安陸府賜并王晃火〔帖木兒〕為食邑}；\eventanchor{SYD-1340-YUAN-197}{壬午，江西行省造螺鈿几榻遺燕鐵木兒，詔賜匠者幣帛各一}；\eventanchor{SYD-1340-YUAN-269}{順元宣撫司亦言：「賊列行營為十六所，乞調兵分道備禦}。

\subsection{1341年：元与高丽、安南及地方政权（年序桥接）}
《元史》、《高丽史》、《安南志略》为这一年留下了可回查的年序入口。

\noindent\textbf{元与高丽、安南及地方政权（年序桥接）。}\eventanchor{SY-1341-CHRONOLOGY-BRIDGE}{【C级年序桥接】保留1341 年的多政权并行检索入口；本条不主张所列政权在当年发生直接互动，具体事件须另以单项材料入账}。

\subsection{1342年：元与高丽、安南及地方政权（年序桥接）、元}
《元史》、《高丽史》、《安南志略》为这一年留下了可回查的年序入口。

\noindent\textbf{元。}\eventanchor{SYD-1342-YUAN-054}{命威順王寬徹不花還鎮湖廣}；\eventanchor{SYD-1342-YUAN-126}{安豐路饑，賑糶麥四萬二千四百石}；\eventanchor{SYD-1342-YUAN-198}{丙申，命參知政事納麟監繪明宗皇帝御容}；\eventanchor{SYD-1342-YUAN-270}{丙子，命興元府鳳州留壩鎮及晉寧路遼山縣十八盤各立巡檢司}。

\noindent\textbf{元与高丽、安南及地方政权（年序桥接）。}\eventanchor{SY-1342-CHRONOLOGY-BRIDGE}{【C级年序桥接】保留1342 年的多政权并行检索入口；本条不主张所列政权在当年发生直接互动，具体事件须另以单项材料入账}。

\subsection{1343年：元与高丽、安南及地方政权（年序桥接）、元}
《元史》、《高丽史》、《安南志略》为这一年留下了可回查的年序入口。

\noindent\textbf{元。}\eventanchor{SYD-1343-YUAN-055}{是月，臨江路新淦州、新喻州，瑞州民饑，賑糶米二萬石}；\eventanchor{SYD-1343-YUAN-127}{以米八千石、鈔二千八百錠，賑哈剌奴兒饑民}；\eventanchor{SYD-1343-YUAN-199}{〔乙〕丑，命宗王燕帖木兒為大宗正府札魯忽赤}；\eventanchor{SYD-1343-YUAN-271}{棒胡本陳州人，名閏兒，以燒香惑眾，妄造妖言作亂，破歸德府鹿邑，焚陳州，屯營於杏岡，命河南行省左丞慶童領兵討之}。

\noindent\textbf{元与高丽、安南及地方政权（年序桥接）。}\eventanchor{SY-1343-CHRONOLOGY-BRIDGE}{【C级年序桥接】保留1343 年的多政权并行检索入口；本条不主张所列政权在当年发生直接互动，具体事件须另以单项材料入账}。

\subsection{1344年：元}
《元史》为这一年留下了可回查的年序入口。

\noindent\textbf{元。}\eventanchor{SYD-1344-YUAN-056}{:洪惟我太祖皇帝混一海宇，爰立定制，以一統緒，宗親各受分地，勿敢妄生覬覦，此不易之成規，萬世所共守者也}；\eventanchor{SYD-1344-YUAN-128}{:於戲，朕豈有意於天下哉！重念祖宗開創之艱，恐隳大業，是以勉徇輿情}；\eventanchor{SYD-1344-YUAN-200}{賜西安王阿剌忒納失里、鎮南王帖木兒不花、威順王寬徹不花、宣靖王買奴等，金各五十兩、銀各五百兩、幣各三十匹}；\eventanchor{SYD-1344-YUAN-272}{帝謂中書省臣曰：「朕在瓊州、建康時，撒迪皆從，備極艱苦，其賜鹽引六萬，俾規利以贍其家}。

\subsection{1345年：元与高丽、安南及地方政权（年序桥接）}
《元史》、《高丽史》、《安南志略》为这一年留下了可回查的年序入口。

\noindent\textbf{元与高丽、安南及地方政权（年序桥接）。}\eventanchor{SY-1345-CHRONOLOGY-BRIDGE}{【C级年序桥接】保留1345 年的多政权并行检索入口；本条不主张所列政权在当年发生直接互动，具体事件须另以单项材料入账}。

\subsection{1346年：元与高丽、安南及地方政权（年序桥接）、元}
《元史》、《高丽史》、《安南志略》为这一年留下了可回查的年序入口。

\noindent\textbf{元。}\eventanchor{SYD-1346-YUAN-057}{丙辰，賜雲南行省參知政事不老三珠虎符，以兵討死可伐}；\eventanchor{SYD-1346-YUAN-129}{丙午，命中書右丞太平、樞密副使姚庸、御史中丞張起巖知經筵事}；\eventanchor{SYD-1346-YUAN-201}{丙戌，開京師金口河，深五十尺，廣一百五十尺，役夫一十萬}；\eventanchor{SYD-1346-YUAN-273}{大名路饑，以鈔萬二千錠賑之}。

\noindent\textbf{元与高丽、安南及地方政权（年序桥接）。}\eventanchor{SY-1346-CHRONOLOGY-BRIDGE}{【C级年序桥接】保留1346 年的多政权并行检索入口；本条不主张所列政权在当年发生直接互动，具体事件须另以单项材料入账}。

\subsection{1347年：元与高丽、安南及地方政权（年序桥接）}
《元史》、《高丽史》、《安南志略》为这一年留下了可回查的年序入口。

\noindent\textbf{元与高丽、安南及地方政权（年序桥接）。}\eventanchor{SY-1347-CHRONOLOGY-BRIDGE}{【C级年序桥接】保留1347 年的多政权并行检索入口；本条不主张所列政权在当年发生直接互动，具体事件须另以单项材料入账}。

\subsection{1349年：元与高丽、安南及地方政权（年序桥接）、元}
《元史》、《高丽史》、《安南志略》为这一年留下了可回查的年序入口。

\noindent\textbf{元。}\eventanchor{SYD-1349-YUAN-058}{甲申，常州宜興山水出，勢高一丈，壞民廬}；\eventanchor{SYD-1349-YUAN-130}{丙戌，加封瀏陽州道吾山龍神崇惠昭應靈顯廣濟侯}；\eventanchor{SYD-1349-YUAN-202}{丙子，開上都、興和等處酒禁}；\eventanchor{SYD-1349-YUAN-274}{丁丑，封皇姊月魯公主為昌國大長公主}。

\noindent\textbf{元与高丽、安南及地方政权（年序桥接）。}\eventanchor{SY-1349-CHRONOLOGY-BRIDGE}{【C级年序桥接】保留1349 年的多政权并行检索入口；本条不主张所列政权在当年发生直接互动，具体事件须另以单项材料入账}。

\subsection{1350年：元与高丽、安南及地方政权（年序桥接）、元}
《元史》、《高丽史》、《安南志略》为这一年留下了可回查的年序入口。

\noindent\textbf{元。}\eventanchor{SYD-1350-YUAN-059}{:每念治必本於盡孝，事莫先於正名，賴天之靈，權奸屏黜，盡孝正名，不容復緩，永惟鞠育罔極之恩，忍忘不共戴天之義}；\eventanchor{SYD-1350-YUAN-131}{:昔我皇祖武宗皇帝昇遐之後，祖母太皇太后惑於憸慝，俾皇考明宗皇帝出封雲南}；\eventanchor{SYD-1350-YUAN-203}{以太保馬札兒台為太師、中書右丞相，太尉塔失海牙為太傅，知樞密院事塔馬赤為太保，御史大夫脫脫為知樞密院事，汪家奴為中書平章政事，嶺北行省平章政事也先帖木兒為御史大夫}；\eventanchor{SYD-1350-YUAN-275}{罷通州、河西務等處抽分按利房，大都東裏山查提領所}。

\noindent\textbf{元与高丽、安南及地方政权（年序桥接）。}\eventanchor{SY-1350-CHRONOLOGY-BRIDGE}{【C级年序桥接】保留1350 年的多政权并行检索入口；本条不主张所列政权在当年发生直接互动，具体事件须另以单项材料入账}。

\subsection{1351年：元}
《元史》为这一年留下了可回查的年序入口。

\noindent\textbf{元。}\eventanchor{SYD-1351-YUAN-060}{戊寅，詔：「作新風憲}；\eventanchor{SYD-1351-YUAN-132}{寶慶路饑，判官文殊奴以所受敕牒貸官糧萬石賑之}；\eventanchor{SYD-1351-YUAN-204}{丁未，立四川省檢校官}；\eventanchor{SYD-1351-YUAN-276}{冬十月乙未，增立巡防捕盜所於永昌}。

\subsection{1352年：元与高丽、安南及地方政权（年序桥接）、元}
《元史》、《高丽史》、《安南志略》为这一年留下了可回查的年序入口。

\noindent\textbf{元。}\eventanchor{SYD-1352-YUAN-061}{庚寅，河決曹州，雇夫萬五千八百修築之}；\eventanchor{SYD-1352-YUAN-133}{丙午，命太平提調都水監}；\eventanchor{SYD-1352-YUAN-205}{丙戌，賜脫脫金十錠、銀五十錠、鈔萬錠、幣帛二百匹，辭不受}；\eventanchor{SYD-1352-YUAN-277}{丁卯，山東霖雨，民饑相食，賑之}。

\noindent\textbf{元与高丽、安南及地方政权（年序桥接）。}\eventanchor{SY-1352-CHRONOLOGY-BRIDGE}{【C级年序桥接】保留1352 年的多政权并行检索入口；本条不主张所列政权在当年发生直接互动，具体事件须另以单项材料入账}。

\subsection{1353年：元与高丽、安南及地方政权（年序桥接）、元}
《元史》、《高丽史》、《安南志略》为这一年留下了可回查的年序入口。

\noindent\textbf{元。}\eventanchor{SYD-1353-YUAN-062}{:朕自踐祚以來，至今十有餘年，託身億兆之上，端居九重之中，耳目所及，豈能周知}；\eventanchor{SYD-1353-YUAN-134}{己卯，監察御史不答失里請罷造作不急之務}；\eventanchor{SYD-1353-YUAN-206}{丙午，命也先帖木兒、鐵木兒塔識並為御史大夫}；\eventanchor{SYD-1353-YUAN-278}{大都、永平、鞏昌、興國、安陸等處并桃溫萬戶府各翼人民饑，賑之}。

\noindent\textbf{元与高丽、安南及地方政权（年序桥接）。}\eventanchor{SY-1353-CHRONOLOGY-BRIDGE}{【C级年序桥接】保留1353 年的多政权并行检索入口；本条不主张所列政权在当年发生直接互动，具体事件须另以单项材料入账}。

\subsection{1354年：元}
《元史》为这一年留下了可回查的年序入口。

\noindent\textbf{元。}\eventanchor{SYD-1354-YUAN-063}{復立八百宣慰司，以土官韓部襲其父爵}；\eventanchor{SYD-1354-YUAN-135}{八月丙午〔朔〕，命江浙行省右丞忽都不花、江西行省右丞禿魯統軍合討羅天麟}；\eventanchor{SYD-1354-YUAN-207}{丙申，以朵兒直班為中書右丞，答兒麻為參知政事}；\eventanchor{SYD-1354-YUAN-279}{丙戌，以遼陽吾者野人等未靖，命太保伯撒里為遼陽行省左丞相鎮之}。

\subsection{1355年：元与高丽、安南及地方政权（年序桥接）、元}
《元史》、《高丽史》、《安南志略》为这一年留下了可回查的年序入口。

\noindent\textbf{元。}\eventanchor{SYD-1355-YUAN-064}{帝牽於眾請，令三年後減之}；\eventanchor{SYD-1355-YUAN-136}{庚辰，詔建木華黎、伯顏祠堂於東平}；\eventanchor{SYD-1355-YUAN-208}{丙戌，亦憐只答兒反，遣兵討之}；\eventanchor{SYD-1355-YUAN-280}{丙戌，中書省臣建議，以河南盜賊出入無常，宜分撥達達軍與揚州舊軍於河南水陸關隘戍守，東至徐、邳，北至夾馬營，遇賊掩捕，從之}。

\noindent\textbf{元与高丽、安南及地方政权（年序桥接）。}\eventanchor{SY-1355-CHRONOLOGY-BRIDGE}{【C级年序桥接】保留1355 年的多政权并行检索入口；本条不主张所列政权在当年发生直接互动，具体事件须另以单项材料入账}。

\subsection{1356年：元}
《元史》为这一年留下了可回查的年序入口。

\noindent\textbf{元。}\eventanchor{SYD-1356-YUAN-065}{平江、松江水災，給海運糧十萬石賑之}；\eventanchor{SYD-1356-YUAN-137}{詔：「京官三品以上，歲舉守令一人，守令到任三月，亦舉一人自代}；\eventanchor{SYD-1356-YUAN-209}{詔翰林國史院纂修后妃、功臣列傳，學士承旨張起巖、學士楊宗瑞、侍講學士黃溍為總裁官，左丞相太平、左丞呂思誠領其事}；\eventanchor{SYD-1356-YUAN-281}{八年春正月戊戌朔，命也先帖木兒知樞密院事}。

\subsection{1357年：元与高丽、安南及地方政权（年序桥接）}
《元史》、《高丽史》、《安南志略》为这一年留下了可回查的年序入口。

\noindent\textbf{元与高丽、安南及地方政权（年序桥接）。}\eventanchor{SY-1357-CHRONOLOGY-BRIDGE}{【C级年序桥接】保留1357 年的多政权并行检索入口；本条不主张所列政权在当年发生直接互动，具体事件须另以单项材料入账}。

\subsection{1358年：元与高丽、安南及地方政权（年序桥接）}
《元史》、《高丽史》、《安南志略》为这一年留下了可回查的年序入口。

\noindent\textbf{元与高丽、安南及地方政权（年序桥接）。}\eventanchor{SY-1358-CHRONOLOGY-BRIDGE}{【C级年序桥接】保留1358 年的多政权并行检索入口；本条不主张所列政权在当年发生直接互动，具体事件须另以单项材料入账}。

\subsection{1359年：元与高丽、安南及地方政权（年序桥接）}
《元史》、《高丽史》、《安南志略》为这一年留下了可回查的年序入口。

\noindent\textbf{元与高丽、安南及地方政权（年序桥接）。}\eventanchor{SY-1359-CHRONOLOGY-BRIDGE}{【C级年序桥接】保留1359 年的多政权并行检索入口；本条不主张所列政权在当年发生直接互动，具体事件须另以单项材料入账}。

\subsection{1360年：元与高丽、安南及地方政权（年序桥接）}
《元史》、《高丽史》、《安南志略》为这一年留下了可回查的年序入口。

\noindent\textbf{元与高丽、安南及地方政权（年序桥接）。}\eventanchor{SY-1360-CHRONOLOGY-BRIDGE}{【C级年序桥接】保留1360 年的多政权并行检索入口；本条不主张所列政权在当年发生直接互动，具体事件须另以单项材料入账}。

\subsection{1361年：元与高丽、安南及地方政权（年序桥接）、元}
《元史》、《高丽史》、《安南志略》为这一年留下了可回查的年序入口。

\noindent\textbf{元。}\eventanchor{SYD-1361-YUAN-066}{〔徭〕賊吳天保復寇沅州}；\eventanchor{SYD-1361-YUAN-138}{八月甲辰，以集賢大學士柏顏為中書平章政事，河南行省平章政事月魯不花為宣政院使}；\eventanchor{SYD-1361-YUAN-210}{丙午，命中書平章政事太不花提調會同館}；\eventanchor{SYD-1361-YUAN-282}{丙寅，命平章政事柏顏提調留守司}。

\noindent\textbf{元与高丽、安南及地方政权（年序桥接）。}\eventanchor{SY-1361-CHRONOLOGY-BRIDGE}{【C级年序桥接】保留1361 年的多政权并行检索入口；本条不主张所列政权在当年发生直接互动，具体事件须另以单项材料入账}。

\subsection{1362年：元与高丽、安南及地方政权（年序桥接）、元}
《元史》、《高丽史》、《安南志略》为这一年留下了可回查的年序入口。

\noindent\textbf{元。}\eventanchor{SYD-1362-YUAN-067}{二月丙戌〔朔〕，詔加封天妃父種德積慶侯，母育聖顯慶夫人}；\eventanchor{SYD-1362-YUAN-139}{庚午，命樞密院以軍士五百修築白河堤}；\eventanchor{SYD-1362-YUAN-211}{己巳，詔天下以中統交鈔壹貫文權銅錢壹千文，準至元寶鈔貳貫，仍鑄至正通寶錢並用，以實鈔法，至元寶鈔通行如故}；\eventanchor{SYD-1362-YUAN-283}{六月壬子，星大如月，入北斗，震聲若雷，三日復還}。

\noindent\textbf{元与高丽、安南及地方政权（年序桥接）。}\eventanchor{SY-1362-CHRONOLOGY-BRIDGE}{【C级年序桥接】保留1362 年的多政权并行检索入口；本条不主张所列政权在当年发生直接互动，具体事件须另以单项材料入账}。

\subsection{1363年：元}
《元史》为这一年留下了可回查的年序入口。

\noindent\textbf{元。}\eventanchor{SYD-1363-YUAN-068}{八月丁丑朔，中興地震}；\eventanchor{SYD-1363-YUAN-140}{丙辰，親策進士八十三人，賜朵烈圖、文允中進士及第，其餘賜出身有差}；\eventanchor{SYD-1363-YUAN-212}{丙戌，蕭縣李二及老彭、趙君用攻陷徐州}；\eventanchor{SYD-1363-YUAN-284}{丙子，命大都至汴梁二十四驛，凡馬一匹助給鈔五錠}。

\subsection{1364年：元与高丽、安南及地方政权（年序桥接）、元}
《元史》、《高丽史》、《安南志略》为这一年留下了可回查的年序入口。

\noindent\textbf{元。}\eventanchor{SYD-1364-YUAN-069}{是月，方國珍復劫其黨下海，入黃巖港，台州路達魯花赤泰不花率官軍與戰，死之}；\eventanchor{SYD-1364-YUAN-141}{置安東、安豐分元帥府}；\eventanchor{SYD-1364-YUAN-213}{中書省臣言：「張理獻言，饒州德興三處，膽水浸鐵，可以成銅，宜即其地各立銅冶場，直隸寶泉提舉司，宜以張理就為銅冶場官}；\eventanchor{SYD-1364-YUAN-285}{安陸賊將俞君正復陷荊門州，知州聶炳死之}。

\noindent\textbf{元与高丽、安南及地方政权（年序桥接）。}\eventanchor{SY-1364-CHRONOLOGY-BRIDGE}{【C级年序桥接】保留1364 年的多政权并行检索入口；本条不主张所列政权在当年发生直接互动，具体事件须另以单项材料入账}。

\subsection{1365年：元与高丽、安南及地方政权（年序桥接）}
《元史》、《高丽史》、《安南志略》为这一年留下了可回查的年序入口。

\noindent\textbf{元与高丽、安南及地方政权（年序桥接）。}\eventanchor{SY-1365-CHRONOLOGY-BRIDGE}{【C级年序桥接】保留1365 年的多政权并行检索入口；本条不主张所列政权在当年发生直接互动，具体事件须另以单项材料入账}。

\subsection{1366年：元与高丽、安南及地方政权（年序桥接）、元}
《元史》、《高丽史》、《安南志略》为这一年留下了可回查的年序入口。

\noindent\textbf{元。}\eventanchor{SYD-1366-YUAN-070}{皇后肅良合氏生日，百官進箋，皇后諭沙藍答里等曰：「自世祖以來，正宮皇后壽日，不曾進箋，近年雖行，不合典故}；\eventanchor{SYD-1366-YUAN-142}{八月戊寅，以李國鳳為中書左丞，陳有定為福建行省平章政事}；\eventanchor{SYD-1366-YUAN-214}{丙申，擴廓帖木兒遣朱珍、盧旺屯兵河中，遣關保、虎林赤合兵渡河，會竹貞、商暠，且約李思齊以攻張良弼}；\eventanchor{SYD-1366-YUAN-286}{丙戌，以方國珍為江浙行省左丞相，弟國瑛、國珉，姪明善，並為江浙行省平章政事}。

\noindent\textbf{元与高丽、安南及地方政权（年序桥接）。}\eventanchor{SY-1366-CHRONOLOGY-BRIDGE}{【C级年序桥接】保留1366 年的多政权并行检索入口；本条不主张所列政权在当年发生直接互动，具体事件须另以单项材料入账}。

\subsection{1367年：元}
《元史》为这一年留下了可回查的年序入口。

\noindent\textbf{元。}\eventanchor{SYD-1367-YUAN-071}{庚戌，貊高殺衞輝守禦官余仁輔、彰德守禦官范國英，引軍至清化，聞懷慶有備，遂還彰德，上疏言：「人臣以尊君為本，以盡忠為心，以愛民為務}；\eventanchor{SYD-1367-YUAN-143}{詔以擴廓帖木兒不遵君命，宜黜其兵權，就命貊高討之}；\eventanchor{SYD-1367-YUAN-215}{愛猷識理達臘計安宗社，累請出師}；\eventanchor{SYD-1367-YUAN-287}{八月丙午，詔命皇太子總天下兵馬，其略曰：「元良重任，職在撫軍，稽古徵今，卓有成憲}。

\subsection{1368年：元}
《元史》为这一年留下了可回查的年序入口。

\noindent\textbf{元。}\eventanchor{SYD-1368-YUAN-072}{丙辰，詔皇太子位下立儀衞司，設指揮二員，給二珠金牌，副指揮二員，一珠金牌}；\eventanchor{SYD-1368-YUAN-144}{丙戌，以武衞所管鹽臺屯田八百頃，除軍見種外，荒閑之地，盡付分司農司}；\eventanchor{SYD-1368-YUAN-216}{賜吳王搠思監金二錠、銀五錠、鈔二千錠、幣帛各九匹}；\eventanchor{SYD-1368-YUAN-288}{答失八都魯克復襄陽、樊城有功，陞四川行省右丞，賜金繫腰一}。
